%==============================================================================
% CONFIGURAÇÃO DO DOCUMENTO
%==============================================================================
\documentclass[a4paper,12pt]{extarticle}  % Classe de documento com tamanho A4 e fonte 12pt

%==============================================================================
% PACOTES E BIBLIOTECAS
%==============================================================================

% --- Configuração de página e layout ---
\usepackage[left=30mm,top=30mm,right=30mm,bottom=30mm]{geometry}  % Margens do documento
\usepackage{etoolbox}     % Necessário para a página de capa
\usepackage{adjustbox}    % Ajustes de caixas
\usepackage{fancyhdr}     % Cabeçalho e rodapé personalizados
\setlength{\headheight}{15.05pt} % Corrige erro do fancyhdr sobre headheight
\usepackage{lastpage}     % Referência à última página

% --- Suporte a idiomas e codificação ---
\usepackage[T1]{fontenc}  % Codificação de fontes
\usepackage[utf8]{inputenc}  % Suporte a caracteres UTF-8
\usepackage[portuguese]{babel}  % Suporte ao idioma português
\usepackage{newunicodechar}
\newunicodechar{₋}{$_-$}

% --- Matemática e símbolos ---
\usepackage{amsmath}      % Funções matemáticas avançadas
\usepackage{amsfonts}     % Fontes matemáticas adicionais
\usepackage{mathtools}    % Extensões para o amsmath
\usepackage{amssymb}      % Símbolos matemáticos adicionais
\usepackage{bm}           % Símbolos matemáticos em negrito
\usepackage{siunitx}      % Formatação de unidades do SI

% --- Figuras e elementos visuais ---
\usepackage{graphicx}     % Inclusão de figuras
\usepackage{subcaption}   % Suporte a subfiguras
\usepackage{xcolor}       % Suporte a cores
\usepackage[table]{xcolor} % Cores em tabelas
\usepackage{float}        % Melhor controle do posicionamento de figuras
\usepackage{tikz}         % Criação de diagramas
\usetikzlibrary{shapes.geometric, arrows}

% --- Tabelas e listas ---
\usepackage{booktabs}     % Melhora a aparência das tabelas
\usepackage{multirow}     % Células que ocupam múltiplas linhas
\usepackage{enumitem,kantlipsum}  % Personalização de listas
\usepackage[usestackEOL]{stackengine}  % Manipulação de texto em pilha

% --- Algoritmos e código ---
\usepackage[ruled,vlined]{algorithm2e}  % Algoritmos em pseudocódigo
\usepackage{listings}     % Inclusão de código-fonte
\usepackage{minted}       % Destaque de sintaxe avançado
\usemintedstyle{emacs}    % Estilo de cores para código
\usepackage{tcolorbox}    % Caixas coloridas para código/exemplos

% --- Títulos e formatação de seções ---
\usepackage{titlesec}     % Personalização de títulos de seções
\usepackage{titletoc}     % Personalização do sumário

% --- Referências e links ---
\usepackage[square,numbers,sort]{natbib}  % Gerenciamento de bibliografia
\usepackage{hyperref}     % Hiperlinks no documento
\usepackage[capitalise]{cleveref}  % Referências cruzadas inteligentes

%==============================================================================
% CONFIGURAÇÕES PERSONALIZADAS
%==============================================================================

% --- Configuração de blocos de exemplo ---
\newtcolorbox[auto counter, number within=section]{exampleblock}[2][]{
    colframe=gray!80!black, 
    colback=white, 
    coltitle=black, 
    title=Exemplo \thetcbcounter: #2, 
    #1, 
    fonttitle=\bfseries, 
    sharp corners=southwest
}

% --- Definição de cores ---
\definecolor{azulTitulos}{RGB}{0, 76, 153}  % Define uma cor azul para títulos

% --- Formatação de títulos de seções ---
\titleformat{\section}
  {\color{azulTitulos}\normalfont\Large\bfseries}  % Seções em azul, grande e negrito
  {\thesection}{1em}{}
\titleformat{\subsection}
  {\color{azulTitulos}\normalfont\large\bfseries}  % Subseções em azul, menor e negrito
  {\thesubsection}{1em}{}

% --- Configuração de hiperlinks ---
\hypersetup{
    colorlinks,  % Ativa links coloridos
    linkcolor={black},  % Links internos em preto
    citecolor={blue!50!black},  % Citações em azul escuro
    urlcolor={blue!80!black}  % URLs em azul mais claro
}

% --- Configuração do cabeçalho e rodapé ---
\pagestyle{fancy}
\fancyhf{}
\fancyhead[L]{\leftmark}
\fancyhead[R]{\thepage}
\fancyfoot[L]{CSRF}
\fancyfoot[C]{Laboratório 1}
\fancyfoot[R]{\thepage\ de \pageref{LastPage}}
\renewcommand{\headrulewidth}{0.4pt}
\renewcommand{\footrulewidth}{0.4pt}

% --- Estilos de página especiais ---
% Estilo para páginas de capítulo e primeira página
\fancypagestyle{plain}{
    \fancyhf{}
    \fancyfoot[C]{Título do Documento}
    \fancyfoot[R]{Página \thepage\ de \pageref{LastPage}}
    \renewcommand{\headrulewidth}{0pt}
    \renewcommand{\footrulewidth}{0.4pt}
}

% Estilo para primeira página (sem cabeçalho nem rodapé)
\fancypagestyle{firstpage}{
    \fancyhf{}
    \renewcommand{\headrulewidth}{0pt}
    \renewcommand{\footrulewidth}{0pt}
}

% --- Configurações gerais ---
\linespread{1}  % Espaçamento de linha simples
\newtheorem{theorem}{Theorem}[section]  % Define ambiente para teoremas
\graphicspath{{img/}}  % Pasta para figuras

%==============================================================================
% CORPO DO DOCUMENTO
%==============================================================================
\begin{document}

%==============================================================================
% PÁGINAS PRÉ-TEXTUAIS
%==============================================================================
% --- Página de capa ---
\pagenumbering{gobble}  % Remove a numeração de páginas inicialmente
\thispagestyle{firstpage}
\newcommand{\bonecoUM}{Aniko Costa}          % Nome do primeiro docente
\newcommand{\bonecoDOIS}{Filipe Moutinho}        % Nome do segundo docente
\newcommand{\docenteTRES}{}                    % Espaço para um terceiro docente (vazio)
\newcommand{\alunoUM}{João Pedro Antunes - \textbf{70380}}      % Nome e número do terceiro aluno
\newcommand{\alunoDOIS}{Júlio Lopes- \textbf{70512}}  % Nome e número do segundo aluno
\newcommand{\judeu}{Marco Romao - \textbf{71348}}      % Nome e número do terceiro aluno
\newcommand{\data}{1º Semestre 2025/2026}      % Período académico 
\newcommand{\cadeira}{Co-design e Sistemas Reconfiguráveis}  % cadeira
\newcommand{\titulo}{Laboratório 1}       % Título do trabalho 
\title{Título} %n mudar                        % Título alternativo (não usado na capa)
\newcommand{\departmento}{Departamento de Engenharia Eletrotécnica e de Computadores}  % Nome do departamento
\newcommand{\curso}{Mestrado em Engenharia Eletrotécnica e de Computadores}  % Nome do curso
% -------------------------------- Front content
\begin{center}\leavevmode                      % Inicia o ambiente centralizado da capa
    \normalfont
    \includegraphics[width=0.75\columnwidth]{img/logo.png}  % Logotipo da instituição
    \vskip 1cm                                 % Espaçamento vertical de 1cm
    \textsc{\Large \departmento}\\[1 cm]       % Nome do departamento em maiúsculas pequenas
    {\large \curso}                            % Nome do curso
    \vskip 1cm                                 % Espaçamento vertical de 1cm
    \rule{\linewidth}{0.2 mm}                  % Linha horizontal decorativa
    {\huge \bfseries \titulo \par}             % Título em negrito e tamanho enorme
    \vskip 1cm                                 % Espaçamento vertical de 1cm
    {\Large \bfseries \cadeira \par}          % Subtítulo em negrito e tamanho grande
           
    \rule{\linewidth}{0.2 mm}\\[1.5 cm]        % Linha horizontal decorativa com espaçamento
     
    \begin{minipage}[t]{0.45\textwidth}        % Coluna esquerda para informações dos alunos
    	\begin{flushleft} \large
    		\emph{Alunos:}\\
    		\ifdefempty{\alunoUM}{}{\alunoUM\\}     % Verifica e exibe o aluno 1 se definido
            \ifdefempty{\alunoDOIS}{}{\alunoDOIS\\} % Verifica e exibe o aluno 2 se definido
            \ifdefempty{\judeu}{}{\judeu\\}         % Verifica e exibe o aluno 3 se definido
    	\end{flushleft}
    \end{minipage}
    \begin{minipage}[t]{0.45\textwidth}        % Coluna direita para informações dos docentes
        \begin{flushright} \large
        	{\emph{Docentes:}}\\
        	\ifdefempty{\bonecoUM}{}{\bonecoUM\\}       % Verifica e exibe o docente 1 se definido
        	\ifdefempty{\bonecoDOIS}{}{\bonecoDOIS\\}   % Verifica e exibe o docente 2 se definido
        	\ifdefempty{\docenteTRES}{}{\docenteTRES\\} % Verifica e exibe o docente 3 se definido
        \end{flushright}
    \end{minipage}
    \vfill                                     % Preenche o espaço vertical restante
    {\normalsize \data\par}                    % Data do semestre acadêmico
\end{center}
\cleardoublepage                               % Finaliza a página e inicia uma nova página (ímpar)  % Inclui a página de capa

% --- Sumário e listas ---
\newpage
\pagenumbering{roman}
\setcounter{page}{1}

\tableofcontents
\newpage

% Listas opcionais (comentadas)
%\listoffigures
%\newpage
%\listoftables
%\newpage
%\listofalgorithms  % Lista de algoritmos em pseudocódigo
%\newpage
%\listoflistings  % Lista de algoritmos em formato de código
%\newpage

%==============================================================================
% CONTEÚDO PRINCIPAL
%==============================================================================
\pagenumbering{arabic}  % Inicia numeração arábica (1, 2, 3...)
\setcounter{page}{1}

% --- Seções do documento ---
\input{sections/introduction}
\clearpage

\input{sections/radiosim}
\clearpage

\input{sections/spicesim}
\clearpage

\input{sections/medicoes}
\clearpage

\input{sections/experiencia_rffe}
\clearpage

\input{sections/conlusao}
\clearpage

%==============================================================================
% PÓS-TEXTUAIS
%==============================================================================
% --- Referências bibliográficas ---
\nocite{*}
\bibliographystyle{unsrtnat}  % Estilo de bibliografia numérico
\bibliography{bibliography}
\clearpage  % Garante que o apêndice comece em uma nova página

% --- Apêndice ---
\appendix
% \input{apendix/appendix}  % Descomente se tiver conteúdo no apêndice

\end{document}